\indent ಮೇಲಿನ ವಾಕ್ಯಗಳಲ್ಲಿ ಕಂಡುಬರುವ ಕ್ರಿಯಾಪದಗಳಿಂದ ಕ್ರಿಯೆ ಹಿಂದೆ ನಡೆದಿರುವುದಾಗಿ ತಿಳಿದು ಬರುವ ಕಾರಣ, ಈ ಕ್ರಿಯಾಪದಗಳಿಗೆ ಭೂತಕಾಲದ ಕ್ರಿಯಾಪದಗಳೆಂದು ಹೆಸರು. ಆದುದರಿಂದ ಮೇಲಿನ ಕ್ರಿಯಾ ಪದಗಳಲ್ಲಿರುವ ಅನು, ಅಳು, ಇತು, ಎ, ಎನು, ಅರು, ಅವು, ಇರಿ, ಎವು— ಎಂಬುವು ಭೂತಕಾಲದ ಆಖ್ಯಾತ ಪ್ರತ್ಯಯಗಳು.\par

\vspace{4pt plus 4pt minus 1pt}\noindent{\kbvsecfont\fontsize{11pt}{14.5pt}\selectfont\color{kbvaccentdark} (3)\hspace{0.45em}ಭವಿಷ್ಯತ್‌ ಕಾಲದ ಕ್ರಿಯಾಪದ ರೂಪಗಳು :}\hspace{0.5em}ಮುಂದೆ ಬರುವ ಕಾಲದ ಕ್ರಿಯೆ ನಡೆಯುವುದನ್ನು ತಿಳಿಸುವ ಕ್ರಿಯಾಪದವು ಭವಿಷ್ಯತ್‌ ಕಾಲದ ಕ್ರಿಯಾಪದವೆನಿಸುವುದು.\par

\indent ಈ ಕೆಳಗೆ ಕೊಟ್ಟಿರುವ ವಾಕ್ಯಗಳಲ್ಲಿ ಕಂಡುಬರುವ ಕ್ರಿಯಾಪದಗಳನ್ನು ಗಮನಿಸಿರಿ.\par

\indent (ಪ್ರ.ಪು.ಏ.).... ಅವನು ಓಡುವನು — ಓಡು + ವ + ಅನು\par

\indent (ಪ್ರ.ಸ್ತ್ರೀ.ಏ.)... ಅವಳು ಓಡುವಳು — ಓಡು + ವ + ಅಳು\par

\indent (ಪ್ರ.ನ.ಏ. ).... ಅದು ಓಡುವುದು — ಓಡು + ವ + ಉದು\par

\indent (ಮ.ವಾ.ಏ.).... ನೀನು ಓಡುವೆ — ಓಡು + ವ + ಎ\par

\indent (ಉ.ವಾ.ಏ.).... ನಾನು ಓಡುವೆನು — ಓಡು + ವ + ಎನು\par

\indent (ಪ್ರ.ವಾ ಬ.).... ಅವರು ಓಡುವರು — ಓಡು + ವ + ಅರು\par

\indent (ಪ್ರ.ನ.ಬ. ).... ಅವು ಓಡುವುವು — ಓಡು + ವ + ಉವು\par

\indent (ಮ.ವಾ.ಬ.)..., ನೀವು ಓಡುವಿರಿ — ಓಡು + ವ + ಇರಿ\par

\indent (ಉ.ವಾ.ಬ).... ನಾವು ಓಡುವೆವು — ಓಡು + ವ + ಎವು\par

\indent ಇಲ್ಲಿ ಒಂಭತ್ತು ಕ್ರಿಯಾಪದಗಳಿವೆ. ಈ ಕ್ರಿಯಾಪದಗಳಲ್ಲಿ “ಓಡು” ಧಾತುವಿನ ಮುಂದೆ ಎಲ್ಲ ಕಡೆಗೂ “ವ” ಎಂಬ ಪ್ರತ್ಯಯ ಬಂದಿದೆ ; ಈ ಪ್ರತ್ಯಯದ ಮುಂದೆ ಏಕವಚನದಲ್ಲಿ ಅನು, ಅಳು, ಉದು, ಎ, ಎನು—ಎಂಬ ಆಖ್ಯಾತ ಪ್ರತ್ಯಯಗಳೂ ಬಹುವಚನದಲ್ಲಿ ಅರು, ಉವು, ಇರಿ, ಎವು— ಎಂಬ ಆಖ್ಯಾತ ಪ್ರತ್ಯಯಗಳೂ ಬಂದಿವೆ.\par

\indent ಮೂಲ ಧಾತುವಿಗೂ ಆಖ್ಯಾತ ಪ್ರತ್ಯಯಕ್ಕೂ ನಡುವೆ ಬಂದಿರುವ “ವ” ಎಂಬುದು ಕಾಲಸೂಚಕ ಪ್ರತ್ಯಯ. ಇದು ಭವಿಷ್ಯತ್‌ ಕಾಲವನ್ನು ಸೂಚಿಸು.\par
